%=====================================================
% Identificação do plano de trabalho
%
% Reúne os dados de cabeçalho exigidos pelo Anexo I das normas de TCC do curso
% de Engenharia Elétrica (DELT/UFPR) e mapeia cada um dos dez itens obrigatórios
% da descrição do projeto para a seção do documento que o atende.
%=====================================================

\chapter*{IDENTIFICAÇÃO DO PLANO DE TRABALHO}

\begin{spacing}{1.15}

\noindent
\begin{tabular}{@{}p{0.26\linewidth}p{0.70\linewidth}@{}}
\toprule
\textbf{Estudante} & André Victor Xavier Pires \\[0.4em]
\textbf{Matrícula} & GRR20212735 \\[0.4em]
\textbf{Curso} & Engenharia Elétrica --- Setor de Tecnologia, Universidade
Federal do Paraná \\[0.4em]
\textbf{Ênfase} & Sistemas eletrônicos embarcados (curso noturno) \\[0.4em]
\textbf{Equipe} & Trabalho individual \\
\midrule
\textbf{Orientador} & Prof.\ Dr.\ Leandro dos Santos Coelho --- Departamento de
Engenharia Elétrica e Programa de Pós-Graduação em Engenharia Elétrica
(PPGEE), UFPR \\[0.4em]
\textbf{Coorientador} & Prof.\ Dr.\ Bruno Pohlot Ricobom --- Departamento de
Engenharia Elétrica, Setor de Tecnologia, UFPR. Área de atuação: instrumentação
eletrônica e medição de campo próximo aplicada a compatibilidade eletromagnética.
Contribuição: concepção das placas de teste e condução da validação experimental
por varredura de campo próximo \\
\midrule
\textbf{Disciplinas} & TCC I (período letivo 2026/2) e TCC II (período letivo
2027/1) \\[0.4em]
\textbf{Título} & \makeatletter\@title\makeatother \\
\bottomrule
\end{tabular}

\vspace{2em}

\noindent
Os itens exigidos para a descrição do projeto, conforme o Anexo I das normas do
curso, são atendidos nas seções relacionadas a seguir.

\vspace{0.8em}

\noindent
\begin{tabular}{@{}p{0.56\linewidth}p{0.40\linewidth}@{}}
\toprule
\textbf{Item do Anexo I} & \textbf{Onde é atendido} \\
\midrule
1) Título & Folha de rosto e Seção~\ref{sec:tema} \\[0.35em]
2) Introdução & Capítulo~\ref{cap:introducao} \\[0.35em]
3) Objetivos & Seção~\ref{sec:objetivos} \\[0.35em]
4) Público alvo & Seção~\ref{sec:publico_alvo} \\[0.35em]
5) Metodologia de desenvolvimento & Capítulo~\ref{cap:metodologia} \\[0.35em]
6) Recursos necessários & Seção~\ref{sec:recursos} \\[0.35em]
7) Resultados fundamentais a serem atingidos & Seção~\ref{sec:resultados_esperados}
\\[0.35em]
8) Contribuição esperada para a ênfase e importância para a formação do autor &
Seção~\ref{sec:contribuicao} \\[0.35em]
9) Cronograma & Seção~\ref{sec:cronograma} \\[0.35em]
10) Bibliografia a ser utilizada & Referências \\
\bottomrule
\end{tabular}

\end{spacing}

\vfill
\cleardoublepage

%=====================================================
